\documentclass{article}
\usepackage[UTF8]{ctex}
\usepackage{amsmath}
\usepackage{amssymb}
\usepackage{graphicx}
\usepackage{geometry}
\geometry{a4paper, margin=2cm}

\newcommand{\problemname}[1]{\begin{center}\Large\bfseries #1\end{center}}
\usepackage{enumitem}

\begin{document}
\problemname{程序设计 T2}

蒜头君第一次网购，在提交订单后，平台赠送他两张张运费险抵用券，可以用来抵扣退换货时的运费。已知每张运费险抵用券的抵扣规则是这样的：

\begin{itemize}
  \item \texttt{A} 券：抵扣一半的运费；

  \item \texttt{B} 券：抵扣 $7$ 元的运费。

\end{itemize}
如果蒜头君在收到商品时不满意，想要退货，而退货时的运费标准是：

\begin{itemize}
  \item 包裹重量不大于 $10$ 公斤时，每公斤 $9$ 元；

  \item 超过 $10$ 公斤后，其超过部分每公斤 $2$ 元；

\end{itemize}
现已知：包裹的重量和蒜头君的抵用券使用情况，请问包裹的运费应为多少元？

\subsection*{输入格式}
输入一行，包含一个整数 $X$（$1\leX\le100$）表示包裹的重量，单位：公斤；字符 $Y$ 表示抵用券的使用情况：

\begin{itemize}
  \item \texttt{A} ：使用了抵扣一半运费的抵用券；

  \item \texttt{B} ：使用了抵扣 $7$ 元运费的抵用券；

  \item \texttt{C} ：没有使用运费抵用券。

\end{itemize}
\subsection*{输出格式}
输出只有一行，包含 $1$ 个数，表示包裹的运费。（结果保留一位小数）

\subsection*{样例输入}
\begin{verbatim}
10 B
\end{verbatim}

\subsection*{样例输出}
\begin{verbatim}
83.0
\end{verbatim}


\end{document}
